


















\chapter{Anamnese und Patientengespräch}
\label{chp:AUU}

\emph{Maintainer: Sebastian Gabriel}

\minitoc

Hier lernen Sie wie Sie die bereits gelernten theoretischen Grundlagen in
der Praxis einsetzen, wie und in welcher Reihenfolge Sie im Einsatz
vorgehen, wie Sie Dringlichkeiten einsch\"atzen und Verdachtsdiagnosen
erarbeiten können und wie Sie dabei dann auch noch den Überblick
behalten können.




%========Beginn Abschnitt Emhofer============================================

\section{{\quotedblbase}Hilfe, ein Patient!{\textquotedblleft} - Was
mach{\textquotesingle} ich jetzt?}

\begin{wrapfigure}{r}{4cm}
\includegraphics[width=\linewidth]{./images/teaser/follow-me-kontur.jpg}
\end{wrapfigure}Gerade zu Beginn der T\"atigkeit als medizinische Fachkraft
fühlt man sich beim Eintreffen am Notfallort durch die große Menge an Sinneseindrücken und
Informationen überlastet. Um Sie sicher durch diese Informationsflut zu
leiten, wollen wir Ihnen in diesem Abschnitt ein paar Anhaltspunkte
pr\"asentieren, welche den Einsatzablauf strukturieren und Ihnen und somit auch
dem Patient Sicherheit geben.

\gLayout{0}{Vorgehen am Notfallort:}{
\begin{itemize}
    \item Bevor es los geht: Szeneüberblick und Selbstschutzmaßnahmen
    \item 1. Schritt: \gEs{Ersteinsch\"atzung des Patienten}
    \item 2. Schritt: \gEs{Erstmaßnahmen}
    \item 3. Schritt: \gEs{Durchatmen, nachdenken, planen}
\end{itemize}
}{
\begin{itemize}
    \item Szeneüberblick und Selbstschutzmaßnahmen
    \item Ersteinsch\"atzung des Patienten
    \item Erstmaßnahmen
    \item Durchatmen, nachdenken, planen
\end{itemize}
}{}{}{}

\dictum[Dickie, in: Samuel Shem: House of God.
\cite{HouseOfGodDe}]{Alles in Ordnung, auch wenn Sie in Panik geraten sind und sich beschissen fühlen.
Ich kenne das. Es ist furchtbar, und es wird nicht das letzte Mal sein.
Aber vergessen Sie nicht, was Sie gesehen haben. Regel Nr. 3: Bei Herzstillstand
zuerst den eigenen Puls fühlen.}

W\"ahrend die ersten zwei Schritte ein gezieltes Vorgehen erfordern, bietet der
dritte Schritt Spielraum für die individuelle Untersuchung und Befragung des
Patienten entsprechend den Symptomen. Im Folgenden stellen wir Ihnen die Schritte
im Einzelnen vor.

% \section{Die Anamnese erz\"ahlt uns die
% Geschichte des Patienten}
\index{Anamnese}

\section{Kommunikationsregeln}
\label{topic:kommunikationsregeln}

\begin{wrapfigure}{r}{4cm}
\includegraphics[width=\linewidth]{./images/teaser/imgp0482.jpg}
\end{wrapfigure}
Der Erfolg des Anamnesegespräches ist für das weitere Vorgehen 
entscheidend. Oft erhält man wichtige Information zum Zustand des Patienten
nicht durch die gemessenen Werte sondern durch die richtige Befragung! Der
Erfolg des Gespräches hängt dabei einerseits von Geschick des Gesprächführenden,
andererseits aber auch von der Bereitschaft des Patienten (oder dessen
Umgebung!) ab. 

Im Folgenden möchten wir Ihnen einige \gE{Grundregeln} und wertvolle Tipps
geben, wie Sie Ihr Patientengespräch so gestalten können, dass Sie die
gewünschten Informationen erhalten und ein Vertrauensverhältnis mit dem
Patienten aufbauen können:

\begin{itemize}
    \item Stellen Sie sich mit Ihrem Namen und Ihrer Funktion vor und - so es
    der Notfall zulässt - geben Sie dem Patienten die Hand!
    \item Begeben Sie sich auf die gleiche Ebene (auf \gE{Augenhöhe}) mit dem
    Patienten!
    \item Hören Sie "`aktiv zu"'! Die Antworten sollen bei Ihnen "`ankommen"'
    und nicht "`verloren gehen"'.
    
    Negativbeispiel: \begin{inparadesc}
                     \item[Sanitäter:] \emph{"`Haben Sie Schmerzen?"'}
                     \item[Patient:] \emph{"`Nein."'}
                     \item[Sanitäter:] \emph{"`Wie stark ist der Schmerz?"'}
                     \end{inparadesc}
    \item Versuchen Sie zu verstehen was
    der Patient \gE{meint}! Wiederholen Sie wichtige Informationen um sicher zu
    gehen dass Sie es richtig verstanden haben.
    \item Verwenden Sie eine leicht verständliche Sprache und vermeiden Sie Fachausdrücke!
    \item Stellen Sie immer nur \gE{eine Frage gleichzeitig}.
    \item Informieren Sie den Patient über Ihr Vorgehen!
    \item Versuchen Sie Fragen des Patienten nach bestem Wissen zu
    beantworten, aber erfinden Sie nichts! Lügen Sie den Patienten nicht an!
    \item Treffen Sie keine Aussagen oder Vorhersagen die ausserhalb Ihres
    Kompetenz- und Zuständigkeitsbereiches liegen.
    \item Im Team sollten nicht alle auf den Patient einreden! Der
    Patient sollte wenn möglich nur einen Ansprechpartner haben!
\end{itemize}

\gMarriedNG{Fragentypen}
{
Grundsätzlich unterscheidet man \gEs{offene} und \gEs{geschlossene} Fragen. 

\subparagraph{Offene Fragen} kann der Patient frei, mit seinen eignen Worten
beantworten. Je nach Patient kann man so mit wenig Aufwand sehr viele
Informationen gewinnen. Das problem dabei kann jedoch sein, dass die Informationen ungeordnet und mitunter
auch in der Situation mäßig wichtig sind. 

\begin{quote}
\emph{"`Was haben Sie für ein Problem?"'}

\emph{"`Können Sie den Schmerz beschreiben?"'}
\end{quote}

\subparagraph{Geschlossene Fragen} sind sehr direkt und speziell, der Patient
kann i.d.R. nur mit \emph{"`ja"'}, \emph{"`nein"'} oder \emph{"`weiss nicht"'} antworten.
Sie machen daher eher nur dann Sinn, wenn man ganz spezielle Dinge bewusst
direkt abfragen kann oder muss. \gCa{Achtung!} Manche Patienten neigen dazu
grundsätzlich "`ja"' zu sagen, ohne dass sie die Frage wirklich verstanden
haben!

\begin{quote}
\emph{"`Haben Sie sich vorher angestrengt?"'}

\emph{"`Ist der Schmerz stechend?"'}
\end{quote}

\subparagraph{Suggestivfragen} sind ganz "`böse"': Hier
legt der Gesprächsführer den Patienten die Antworten "`in den Mund"'. 

\begin{quote}
\emph{"`Haben Sie \emph{eh keine} Schmerzen?"'}

\emph{"`Ist der Schmerz stechend?"'}\footnote{Ja, das ist die gleiche Frage wie
das Beispiel bei der geschlossenen Frage. Man sieht wie schnell es gehen kann,
ohne es zu wllen eine Suggestivfrage zu stellen \ldots}
\end{quote}

Dieser Fragentyp
soll nicht -- oder nur im begründeten Ausnahmefall -- angewendet werden. 

}
{
% TODO VIT: Slide fehlt!
\begin{itemize}
    \item Offene Fragen.
    \item Geschlossene Fragem.
    \item Suggestivfragen (böse!).
\end{itemize}
}

\gMarriedNG{Kinder}
{ Bei Notfällen mit Kindern kann die Kommunikation besonders herausfordernd
werden, da Kinder, die Angst haben, sich oft gegenüber Fremden verschließen. Die
oben genannten Regeln gelten daher bei Kindern umso mehr. Zusätzlich können
folgende Tipps hilfreich sein:

\begin{itemize}
    \item Kinder nehmen non-verbale Signale stärker wahr als Erwachsene. Versuchen Sie daher durch einen freundlichen Gesichtsausdruck und eine offene Körperhaltung Vertrauen zu bilden!
    \item Erklären Sie dem Kind (altersgerecht) warum Sie gerufen wurden!
    \item Sprechen Sie das Kind direkt und mit seinem Namen an! Vermeiden Sie es über das Kind in der 3. Person zu sprechen, damit das Kind nicht das Gefühl bekommt, dass hinter seinem Rücken getuschelt wird. Dies würde dem Kind unnötig Angst machen!
    \item Seien Sie immer ehrlich! Kinder merken sehr schnell, wenn Sie lügen. Lügen Sie ein Kind an, so wird das Kind schnell das Vertrauen in Sie verlieren und jede weitere Untersuchung verweigern!
    \item Geben Sie dem Kind ein Spielzeug in die Hand!
    \item Trennen Sie das Kind niemals von der Mutter!
    \item Hat das Kind Angst vor einer bestimmten Untersuchung oder einer Maßnahme, zeigen Sie diese zuerst an der Mutter oder an einem Stofftier vor!
    \item Loben Sie das Kind nach der Untersuchung bzw. nach der Maßnahme!
    \item Da Kinder ihre Beschwerden nicht eindeutig beschreiben können, lassen Sie sich z.B. den Ort der Schmerzen zeigen.
\end{itemize}
}{
{\gSymbolText}
}


\section{Was muss ich erfragen? -- S-KAMEL}

\index{S-KAMEL}

Neben der körperlichen Untersuchung gibt das Gespr\"ach mit dem Patienten
wertvolle Hinweise auf das Notfallgeschehen. Diese Erhebung der Kranken- bzw.
Unfallgeschichte nennt man \gT{Anamnese}. Eine gute Anamnese ist mindestens
genauso wichtig wie die Erhebung der Befunde! Leider wird in der Praxis oft der
Messwerterhebung eine größere Bedeutung beigemessen.

\gLayout{0}{S-KAMEL}{
Auch das Anamnesegespräch soll strukturiert durchgeführt werden. Dafür stehen
mehrere Abfrageschemata zur Verfügungen, von welchen wir Ihnen ein in Österreich
weit verbreitetes vorstellen möchten: Das \gT{S-KAMEL-Schema}. Nach einer
Einstiegsfrage, z.B. "`Was können wir für Sie tun?"', beinhaltet das
S-KAMEL-Schema folgende Punkte:

\begin{itemize}
    \item Symptome und Schmerzen
    \item Krankheiten
    \item Allergien
    \item Medikamente
    \item Ereignisse vor Notfalleintritt
    \item Letzte...
\end{itemize}

Dieses Abfrageschema kann Ihnen im hektischen Notfall
Sicherheit geben. Durch das Abarbeiten der Fragen haben Sie
die Sicherheit, dass keine wichtigen Informationen vergessen
werden. Außerdem ist es für den Patienten vertrauensbildend,
da die gezielte Befragung professionell wirkt. Im Folgenden
müssen wir uns die einzelnen Punkte des Schemas genauer
ansehen.
}{
\begin{itemize}
    \item Symptome und Schmerzen
    \item Krankheiten
    \item Allergien
    \item Medikamente
    \item Ereignisse vor Notfalleintritt
    \item Letzte...
\end{itemize}
}{}{}{}

\gFazit{S{\textquotesingle}KAMEL denkt an alles.}

\subsubsection{Symptome und Schmerzen}
\index{Schmerz}

Die Einstiegsfrage hat uns u.U. bereits auf das Hauptsymptom hingewiesen. Ist dies nicht der Fall, so muss als erstes nach dem Hauptsymptom gefragt werden. Fragen können z.B. sein:
\begin{itemize}
    \item \emph{"`Welche Beschwerden haben Sie?"'}
    \item \emph{"`Wo drückt der Schuh?"'}
\end{itemize}
Außerdem muss immer auch speziell nach Schmerzen gefragt werden!
\begin{itemize}
    \item \emph{"`Haben Sie Schmerzen?"'}
\end{itemize}

\gLayout{0}{Fragen über Schmerzen:}{
Gibt der Patient an Schmerzen zu haben, müssen wir weitere Informationen
erfragen, damit wir uns ein klares Bild über die Beschwerden machen können:
\index{Schmerz!Arten}

\begin{itemize}
    \item \gEs{Lokalisation}: \emph{"`Wo sind die
    Schmerzen?"'}
    \item \gEs{Zeitdauer/Beginn}: \emph{"`Seit wann haben Sie
    die Schmerzen? Wie lange haben Sie die Schmerzen schon?
    Haben Sie so einen Schmerz schon einmal gehabt? Haben die
    Schmerzen plötzlich oder langsam eingesetzt?"'} 
    \item
    \gEs{Qualität}: \emph{"`Wie ist der Schmerz? Bohrend,
    stechend, ziehend, krampfartig, auf- und
    abschwellend?"'}
    \item \gEs{Ausstrahlung}: \emph{"`Strahlt der Schmerz
    irgendwohin aus?"'} 
    \item \gEs{Provokation}: \emph{"`Was
    macht den Schmerz besser oder schlechter?"'} 
    \item \gEs{Stärke}: \emph{"`Wie stark ist der
    Schmerz?"'} \gZ{In einigen Lehrbüchern wird empfohlen, den Schmerz anhand
    einer Skala von 1-10 beurteilen zu lassen, wir finden
    jedoch, dass auch eine solche Beurteilung recht subjektiv
    ist und daher nur mit Vorsicht verwendet werden sollte.}
    Setzen Sie die Aussage des Patienten immer in Relation zu
    seinem Gesichtsausdruck und zu seiner Körperhaltung!
    Einem Patient mit starken Schmerzen sieht man dies an!
\end{itemize}
}{
\begin{itemize}
  \item Wo?
  \item Ausstrahlung?
  \item Seit wann? Wie lange schon?
  \item Plötzlicher oder langsamer Beginn?
  \item Verlauf?
  \item Wie?
  \item Wie stark?
  \item Was provoziert den Schmerz?
\end{itemize}
}{}{}{}

\paragraph{Besonderheit: \gT{Kolikartiger Schmerz}}~
\label{topic:kolik}
\index{Kolik!Schmerzqualität}
\index{Schmerz!kolikartiger}

\gDefinition{Ein kolikartiger Schmerz ist auf-
und abschwellend.}

Er wird oft als krampfhaft und sehr stark beschrieben. Er
kommt häufig bei Verstopfung der Gallen- und Harnwege durch
Steine vor (\gSee{topic:gallenkolik},
\ref{topic:nierenkolik}).

\subsubsection{Krankheiten}
Notfälle ereignen sich manchmal aufgrund der akuten Verschlechterung einer bestehenden Vorerkrankung. Außerdem hilft die Beurteilung der Messwerte, wenn man über die Grunderkrankungen des Patienten Bescheid weiß. Zum Beispiel kann die Ursache eines hohen Blutdrucks eine bereits bestehende Erkrankung aber auch ein Notfallgeschehen sein. Neben der Befragung kann auch ein alter Spitalsbericht (Arztbrief) Auskunft über bestehende Vorerkrankungen liefern. Allgemeine Fragen, wie z.B.
\begin{itemize}
    \item \emph{"`Leiden Sie unter irgendwelchen
    Krankheiten?}
    \item \emph{"`Bestehen bei Ihnen Vorerkrankungen?}
\end{itemize}
können anschließend durch symptom-spezifische Fragen ergänzt werden:
\begin{itemize}
    \item \emph{"`Hatten Sie schon einmal einen
    Herzinfarkt?"'}
    \item \emph{"`Leiden Sie unter hohem Blutdruck?"'}
    \item \emph{"`Sind Sie zuckerkrank?"'}
\end{itemize}

\subsubsection{Allergien}
Auf die Frage nach möglichen Allergien wird häufig vergessen, ein Informationsmangel kann aber fatale Auswirkungen haben. Fragen Sie daher immer nach Allergien, in einer zweiten Frage sogar speziell nach Medikamentenallergien! Damit können Sie für Not- oder Spitalsarzt wichtige Informationen erheben. (Natürlich könnten das die Ärzte auch selbst fragen, aber denken Sie daran, dass ein Notfallpatient bewusstlos werden kann und es deshalb wichtig ist, dass sie diese Informationen einholen, solange der Patient noch ansprechbar ist!)

\subsubsection{Medikamente}
Die vom Patienten eingenommenen Medikamente lassen
Rückschlüsse auf seine Grunderkrankungen zu und sind im Falle
eines längeren Spitalsaufenthalts eine wertvolle Information.
Auch die Regelmäßigkeit oder Unregelmäßigkeit der Einnahme
kann Ihnen Anhaltspunkte für das Finden Ihrer
Verdachtsdiagnose bzw. der Umstände des Patienten liefern.

\gLayout{0}{Fragen zu Medikamenten:}{
\begin{itemize}
    \item \emph{"`Welche Medikamente nehmen Sie?"'}
    \item \emph{"`Wogegen nehmen Sie diese Medikamente?"'}
    \item \emph{"`Nehmen Sie Ihre Medikamente regelmäßig?"'}
    \item \emph{"`Wann haben Sie Ihre Medikamente zuletzt"'}
    genommen?
\end{itemize}
}{
\begin{itemize}
  \item Welche?
  \item Wogegen?
  \item Regelmäßige Einnahme?
  \item Wann zuletzt genommen?
\end{itemize}
}{}{}{}

\subsubsection{Ereignisse vor Notfalleintritt}
Oft ereignen sich vor einem Notfall Schlüsselereignisse, welche den Notfall herbei führen oder diesen auslösen. Fragen Sie daher immer nach möglichen Besonderheiten oder Tätigkeiten des Patienten vor dem Notfall:
\begin{itemize}
    \item \emph{"`Was haben Sie vor Beginn der Symptome
    gemacht?"'}
    \item \emph{"`Haben Sie sich vor Beginn der Symptome
    besonders angestrengt?"'}
    \item u.v.m.
\end{itemize}
In dieser Kategorie kann es auch hilfreich sein, wenn Sie Angehörige oder Zeugen befragen:
\begin{itemize}
    \item Hat der Patient vor dem Notfall etwas Eingenommen (Überdosis)?
    \item Hat der Patient gekrampft bevor er bewusstlos liegen geblieben ist?
    \item Ist der Patient vor dem Eintreten der Symptome gestürzt?
    \item u.v.m.
\end{itemize}

\subsubsection{Letzte ...}
Der letzte Punkt des Abfrageschemas bezieht sich auf letzte
Ereignisse, welche mit dem Notfall in Zusammenhang stehen
könnten. Damit Sie hier die richtigen Fragen stellen können,
ist ein fundiertes Fachwissen notwendig. Abhängig von den
Beschwerden können verschiedene Fragen zielführend sein.

\gLayout{0}{Fragen zu Letzte\ldots}{
\begin{itemize}
    \item \emph{"`Wann haben Sie zuletzt etwas
    gegessen/getrunken?"'} (Ist der Patient nüchtern?)
    \item \emph{"`Wann hatten Sie zuletzt Stuhlgang?"'}
    (weiterführende Frage: Wie hat der Stuhl ausgesehen?)
    \item \emph{"`Wann war die letzte Regelblutung?"'}
    (Besteht die Möglichkeit einer Schwangerschaft?)
    \item \emph{"`Wann waren Sie zuletzt beim Arzt bzw. im
    Spital?"'}
\end{itemize}
}{
\begin{itemize}
  \item Letzte Mahlzeit?
  \item Letzter Stuhlgang?
  \item Letzte Regelblutung?
  \item Letzter Spitalsaufenthalt?
  \item Letzter Arztbesuch?
  \item Letzte Blutzuckerkontrolle?
  \item \ldots
\end{itemize}
}{}{}{}

\section{Der Arztbrief -- Ein besonderes Hilfsmittel für
die Anamnese}

\subsubsection{Beispiel: Patient Zapfl}







{\ttfamily




\begin{flushright}
 Krankenanstalt der Barmherzigkeit \\II. Med. Abteilung
 
 Hartlgasse 16a \\A-1240 Wien
 
+43\,555\,328\,745\,0\\KADBMED2@grissu.org

\end{flushright}

                                                                                               

\vspace{1.5em}
                                                                                                                            


 Peter Zapfl\par
Diesseitsweg 13a/12/5/34 \par
A-1147 Wien
\vspace{3em}

\begin{flushright}
Wien, am 25.07.2009
\end{flushright}

  Wir berichten über den Aufenthalt unseres Patienten Herrn Peter ZAPFL, geb.
  am 13.12.1936, welcher vom 30.03.2009 bis 17.04. 2009 an unserer Abteilung
  stationär aufgenommen war.

{\bfseries   Aus der Anamnese:}

\begin{quote}
Kindheit: \gMarginNo{\gZ{AE: Entfernung des Blinddarmes; TE: 
Entfernung der Mandeln.}}AE, sowie TE, WK II: längerer Lazarettaufenthalt wg.
Typhus.

Seit 2003 sind ein bis dato nicht insulinpflichtiger Diabetes mellitus, sowie
eine art. Hypertonie bekannt. 

2005 Aufenthalt an unserer Abteilung wg. eines im \gMarginNo{\gZ{CT:
Computertomographie.}}CT nachgewiesenen rechtshirnigen ischämischen Insultes,
anschl. Rehabaufenthalt in Bad Schallerbach.

  Zuletzt lag Pat vom 21.11.2008 bis 23.12.2008 wegen eines nicht
  insulinpflichtigen, entgleisten Diabetes mell., bei Z. n ischäm. Zerebralem.
  Insult 2005 mit \gMarginNo{\gZ{HSZ: Halbseitenzeichen.}}HSZ li und Art.
  Hypertonie an unserer Abteilung.

Zwischenzeitlich keine Auffälligkeiten.

Alkohol.: verneint. Nikotin: verneint. Koffein: fallweise. Harn, Stuhl: unauff.
Schlaf: unauff. Ven.: neg
\end{quote}


{\bfseries Bisherige Med.:} 

\begin{quote}
\begin{tabular}{lll}
Renitec 10 mg &&1-1-1-
\\
Diabetex 850 mg &&1-1-1 
\\
Diab Diät 12 BE &&2-1-4-2-2-1
\\
Thromboass &&1-0-0
\\
\end{tabular}
\end{quote}





 Die nunmehrige Aufnahme erfolgte mit dem ASB, der wegen zunehmender
 Verschlechterung des AZ des Pat. vom Wohnungsnachbarn berufen wurde.

{\bfseries Auffälliges im Aufnahmestatus:}

\begin{quote}
 Deutlich herabgesetzter \gMarginNo{AZ: Allgemeinzustand\gZ{, EZ:
 Ernährungszustand}.}AZ, noch ausreichender EZ, zeitlich und örtlich
 teilorientiert, motorische Unruhe. Aktive Beweglichkeit der re.
 \gMarginNo{OE: Obere Extremität, UE: Untere Extremität.}OE und UE, beinbetonte
 spastische Halbseitenlähmung li mit gesteigerten Sehnenreflexen li bei Z.n. re-Hirn-Insult, mäßige Artikulationsstörungen.
 
 Haut trocken, \gMarginNo{\gZ{Turgor: vom Flüssigkeitsgehalt abhängiger
 Spannungszustand der Haut. \nocite{Pschy259}}}Turgor herabgesetzt, kein
 Dekubitus nachweisbar.
 \end{quote}

{\bfseries Befunde:}
\begin{quote}
Labor siehe Ablichtung.

\gMargin{\gZ{C/P: "`Cor-Pulmo"', Thoraxröntgen.}}C/P.: Keine wesentl Änderung
gegenüber Vorbefund 12.2008

Schädel CT: Keine wesentliche Änderung gegenüber Letztbefund vom 12.2008,
Feinbefund siehe siehe Ablichtung

Augen:. Erstgradige hypertone Veränderungen der
\gMarginNo{\gZ{Retina: Netzhaut.}}Retinagefässe, kein HW für
\gMarginNo{\gZ{Diabetische Retinopathie: Schädigung der Netzhautgefäße
im Rahmen des Diabetes mellitus.}}diabet. Retinopathie.

Neuro: Anamnst Z. n. Insult 2005 im Strombereich der A. cerebri ant., mit
beinbetonter spastischer Halbseitenlähmung li und gesteigerten
\gMarginNo{\gZ{BSR, TSR: Sehnenreflexe des Bizeps und Trizeps.}}BSR und TSR li
sowie stark gesteigerten \gMarginNo{\gZ{PSR: Patellarsehnenreflex.
Babinski-Reflex: krankhafter Reflex bei Schädigung des Rückenmarks oder
Gehirns.}}PSR li, Babinski li pos, Trömner u. Knips bds neg, kein Hinweis auf
Re-Insult, weiterhin physi\-ka\-lisch- thera\-peutische und logopädische
Maßnahmen empf. Der durchgeführte Schlucktest ergab keinen Hinweis auf eine relevante Schluckstörung.

Uro: Prost. rect. dig. Unauff.

EKG: \gMarginNo{\gZ{Sinusrhythmus, HF 74, Linkstyp.}}SR 74, LT,
altersentsprechender Kurvenverlauf.

RR bei mehrfachen Kontrollen gegen Ende des Aufenthaltes normoton, zuvor im
hypotonen Bereich.
\end{quote}




{\bfseries Epikrise:}

\begin{quote}
  Nach Behebung des bei Aufnahme bestehenden Flüssigkeitsdefizites mit
  geeigneter Infusionstherapie, besserte sich der AZ des Pat. zusehends. Mit in der Folge
  durchgeführten physi\-ka\-lisch- therap. Maßnahmen konnte der
  Mobilisisationsgrad des Pat. günstig beeinflusst werden.
  
Der Pat. wird mit bis zum selbständigen Gehen mit Rollator mobilisiert entlassen.
Die zu Aufenthaltsbeginn durchwegs im hypotonen Bereich gemessenen RR-Werte haben
sich in mehrfachen Kontrollen unter dem angegeben antihypertonen Regime
normalisiert. Ebenso hat sich die ursprünglich unbefriedigend gefundene
diabetische Stoffwechsellage, bei konsequenter Einhaltung des vorgeschlagenen
Regimes im therapeutischen Bereich stabilisiert.

  Im Einverständnis mit dem Pat. und dessen Tochter wurde ein Platz im
  Bernhard-Prischl-Heim zur Gewährleistung einer regelmäßigen und auch
  ausreichenden Flüssigkeitszufuhr, sowie einer gesicherten Diabetes- Diät,
  organisiert. Eine Kontrolle der vor etwa einem Jahr durchgeführten 
  \gMarginNo{\gZ{Carotissonographie: Ultraschall der
  Halsschalgader.}}Carotissonographie ist angezeigt. Ein entspr. Termin wurde am
  27.12.2009  12,30 Uhr an der Angiolog Ambulanz d. KA Mitschka-Stiftung
  fixiert. Eine Befund\-be\-sprech\-ung ist nach Terminvereinbarung an h.o. Int.
  Ambulanz möglich.
\end{quote}

\gMarginNo{\gZ{}}
{\bfseries Hauptdiagnose:}

\begin{quote}
\begin{tabular}{ll}
Hypohydratation & E86.0
\\
\end{tabular}
\end{quote}

{\bfseries Nebendiagnosen:}

\begin{quote}
\begin{tabular}{ll}
entgleister. Diabetes mell (NIDM) & E12.31, E12.40, E12.72 
\\
Art. Hypertonie & I10.1
\\
Z.n. Zerebralem Insult mit HSS li 2005& I63.3
\\
Spastische Hemiparese li, beinbetont&I69.3, G81.1
\\
\end{tabular}
\end{quote}



{\bfseries Therapie:}

\begin{quote}
\begin{tabular}{lll}
Renistad Tab. 5,0 mg && 1-1-1 
\\
Glucophage 850 mg Tab && 1-1-1 
\\
Diab. Diät 12 BE&&(2-1-4-2-2-1) 
\\
Thromboass Tab.&&1-0-0 
\\
tgl. Flktszufuhr nicht unter 2 l&&
\end{tabular}
\end{quote}




{\bfseries Kontrollen:}

\begin{quote}
 NFP, Harnbefund, Elektrolyte, RR, BZ, (Pat. führt seit
Jahren Selbstkontrollen durch und führt darüber Aufzeichnungen)) weiterhin neurolog. und
ophthalmolog. Ko
\end{quote}









Mit kollegialen Grüßen\vspace{2em}\\ Dr. Adolf Schrammel, Stationsarzt

\gH{Schrammel}

vidit OA Univ.Doz. Dr. Alois Schremser
}

\vspace{2em}

\hrule












\chapter{Patientenmanagement}
% \clearpage
\section{Vorgehen in der Praxis}

\subsection{Primäreinsatz}

%\paragraph{Szeneüberblick}
% \paragraph{Die ersten  30 Sekunden }

% \gFazit{Ersteinsch\"atzung ${\rightarrow}$~Erstmaßnahmen
% ${\rightarrow}$~Plan zurechtlegen ${\rightarrow}$~Weitere Untersuchung,
% Anamnese und Versorgung}

\begin{gWideParEnv}
\gBox{Ablauf}{Patientenmanagement}{Beige}
{\sffamily
\begin{tabularx}{\linewidth}{cX}
\gTabHeading{Ersteinschätzung}
&
\noindent\fcolorbox{green!10}{green!10}{
% \begin{minipage}{\linewidth}
\begin{tabular}{ll}
\toprule\addlinespace
\gTabSub{}
&
\gTabSub{Szeneüberblick}
\\\addlinespace\midrule\addlinespace
\gTabSub{\gZ{GI}}
&
 \gTabSub{Ersteindruck}
\\\addlinespace
\gTabSub{\gZ{}}
&
 \gTabSub{Bewusstsein}
\\\addlinespace
\gTabSub{\gZ{}}
&
\gTabSub{Hauptbeschwerde}
\\\addlinespace\midrule\addlinespace
\gTabSub{\gZ{A}}
&
 \gTabSub{Atemweg}
\\\addlinespace
\gTabSub{\gZ{B}}
&
\gTabSub{Atmung}
\\\addlinespace
\gTabSub{\gZ{C}}
&
 \gTabSub{Kreislauf}
\gTabSub{\gZ{}}
\\\addlinespace
\gTabSub{\gZ{D}}
&
\gTabSub{Schnelle Untersuchung}
\\\addlinespace\bottomrule
\end{tabular}
% \end{minipage}
}
\\\addlinespace
${\downarrow}$
&
\\\addlinespace
\noindent\gTabHeading{Erstmaßnahmen}
&
Je nach Ergebnis der Ersteinschätzung. 

Z.B. bei vitaler Bedrohung: \gTxMassVit.
\\\addlinespace
${\downarrow}$
&
\\\addlinespace
\noindent\gTabHeading{Plan zurechtlegen} 
&
\\\addlinespace
${\downarrow}$
&
\\\addlinespace
\noindent\gTabHeading{Weitere Untersuchung,
Anamnese und Versorgung}
&
Übersicht Tafel \gSee{tab:checkliste-patientenmanagement}
\\\addlinespace
\end{tabularx}




% \sffamily 
% \begin{center}
% Ersteinsch\"atzung 
% 
% ${\downarrow}$
% 
% Erstmaßnahmen
% 
% ${\downarrow}$
% 
% Plan zurechtlegen 
% 
% ${\downarrow}$
% 
% Weitere Untersuchung,
% Anamnese und Versorgung
% \end{center}
}
\end{gWideParEnv}



\subsubsection{Erster Schritt: Ersteinsch\"atzung}

\begin{gWideParEnv}
\gMassnahmenNG{Standardmaßnahmen}{Ersteinschätzungsblock}{\label{m:ersteinschaetzung}}{
\noindent\begin{minipage}[t]{\linewidth -
{2\fboxrule} - {2\fboxsep}}
    	\noindent\fcolorbox{green!10}{green!10}{\begin{minipage}[t]{\linewidth}
    	{\vspace{-0.8em} 
% \begin{tabulary}{\linewidth}[H]{cp{8em}p{8.5cm}}
\begin{tabulary}{\linewidth}[H]{llL}
\toprule\addlinespace
\gTabSub{}
&
\gTabSub{Szeneüberblick}
&
Sicherheit, Selbstschutz; Großschaden?; Umgebung, Ort, Zeit; Umstände,
Hinweise auf Unfallmechanismus
\\\addlinespace\midrule\addlinespace
\gTabSub{\gZ{GI}}
&
 \gTabSub{Ersteindruck}
 &
Aussehen, Gesichtsausdruck, Haltung, Sprache

Bl\"asse, Zyanose, Schweiß, Schohnhaltung, Atemnot, verwaschene
Sprache, Des\-orientiertheit, Kr\"ampfe 
\\\addlinespace
\gTabSub{\gZ{}}
&
 \gTabSub{Bewusstsein}
 &
Reaktion auf verbale Ansprache, Weckversuch, Schmerzreiz

klar, somnolent, soporös, komatös/bewusstlos
\\\addlinespace
\gTabSub{\gZ{}}
&
\gTabSub{Hauptbeschwerde}
&
Grund der Berufung, Schmerzen und andere für den Patienten dringliche
Symptome
\\\addlinespace\midrule\addlinespace
\gTabSub{\gZ{A}}
&
 \gTabSub{Atemweg}
 &
Atemger\"ausch, Gestik

frei oder verlegt
\\\addlinespace
\gTabSub{\gZ{B}}
&
\gTabSub{Atmung}
&
Frequenz, Atemmuster, Anstrengung beim Atmen, Atemger\"ausche,
Heimsauerstoff
\\\addlinespace
\gTabSub{\gZ{C}}
&
 \gTabSub{Kreislauf}
\gTabSub{\gZ{}}
&
Bewusstsein, Hautfarbe, Radialispuls, Rekap-Zeit
\\\addlinespace
\gTabSub{\gZ{D}}
&
\gTabSub{Schnelle Untersuchung}
&
\begin{minipage}[t]{\linewidth}
\protect\gParallel{
\gEs{Bei Traumapatienten:} 
\begin{itemize}
    \item 1. Helfer: Vitalwerte 
    \item 2. Helfer: Schneller Traumacheck
% \item Suche nach: Schwere Verletzungen, Becken-, Oberschenkel-,
% Serienrippenfraktur, massive Blutung, hartes Abdomen.
\end{itemize}
}{
\gEs{Bei Nicht-Traumapatienten:} 

\begin{itemize}
    \item Vitalwerte.
\end{itemize}
}
\end{minipage}
\\\addlinespace\bottomrule
\end{tabulary}

    	}
      		\end{minipage}}
		\end{minipage}

}
\end{gWideParEnv}





Nun stellen sich folgende Fragen:

\begin{enumerate}
    \item \emph{Was ist das Problem des Patienten?}
    \item \emph{Ist der Patient akut vital bedroht?} (Alarmzeichen, s.u.)
    \item \emph{\ldots oder sind noch weitere Untersuchungen zur Lagebeurteilung
    nötig? (z.B. Vitalwerte)}
    \item \emph{Welche weiteren Untersuchungen, Maßnahmen und Fragen stehen an
    erster Stelle?}
\end{enumerate}

    	

\gLayout{28}{Vitale Bedrohung -- Alarmzeichen
{\color{ubuntured}\bfseries{\VarFlag ~}}} {
Das rasche Festellen einer vitalen Bedrohung ist besonders wichtig. Anzeichen
einer vitalen Bedrohung können u.a. sein: Störungen von Be\-wusstsein, Atemweg, Atmung oder Kreislauf sowie typische
Symptome wie sich nicht bessernde Atemnot, brodelndes
Atem\-ge\-r\"ausch, Thoraxschmerz, Schocksymptome und entgleiste
Vital\-werte. Manche Diagnosen sind automatisch mit einer vitalen Bedrohung
gleichzusetzen, auch wenn eventuell Symptome fehlen oder nur schwach ausgeprägt
sind. 

\gEs{Sobald eine vitale Bedrohung erkannt wird müssen die \gTxMassVit{}
 (\prettyref{tab:standardmassnahmen-vital}) frühzeitig
ergriffen werden (evt. sogar noch vor Abschluß des Ersteinschätzungsblockes)!}}{\vspace{1em}

\fcolorbox{ubuntured}{White}{
\begin{minipage}[t]{\linewidth - 4\gKorrektur}
{
\begin{itemize}
    \item Massive Störungen von Vitalfunktionen
    \item Typische Symptome
% \begin{itemize}
%     \item Atemnot, die sich nicht bessert 
%     \item Brodelndes Atemger\"ausch 
%     \item Thoraxschmerz  
%     \item Schocksymptome 
%     \item entgleiste Vitalwerte
%     \item schwere Verletzungen
%     \item {\dots} 
% \end{itemize}
\item Gefährliche Diagnosen
\end{itemize}

Übersicht Tafel \gSee{tab:alarmzeichen}

% \begin{itemize}
%     \item Bewusstsein
%     \item Atemweg / Atmung 
%     \item Kreislauf 
% \end{itemize}




}
\end{minipage}
}
}{./images/teaser/minen_hinweis2.jpg}{}{}


% \marginpar{~\newline}
% \gParallel{
% \fcolorbox{ubuntured}{White}{
%       		\begin{minipage}[t]{\linewidth - 2\gKorrektur}
%               \gCa{Alarmzeichen}
%               
%       			{\sffamily\gEs{Massive Störungen von} 
% 
% Bewusstsein
% 
% Atemweg / Atmung 
% 
% Kreislauf 
% 
% \gEs{Typische Symptome} 
% 
% Atemnot, die sich nicht bessert 
% 
% Brodelndes Atemger\"ausch 
% 
% Thoraxschmerz  
% 
% Schocksymptome 
% 
% entgleiste Vitalwerte
% 
% schwere Verletzungen
% 
% {\dots} }
% \end{minipage}
% }}
% {
% \paragraph{Alarmzeichen}
% 
% Anzeichen einer vitalen Bedrohung können u.a. sein: Störungen
%         von Be\-wusstsein, Atemweg, Atmung oder Kreislauf sowie typische
%         Symptome wie sich nicht bessernde Atemnot, brodelndes
%         Atem\-ge\-r\"ausch, Thoraxschmerz, Schocksymptome und entgleiste
%         Vital\-werte.
%         
% \begin{center}
% \includegraphics[width=4cm]{./images/teaser/minen_hinweis2.jpg}
% \end{center}
% }

\gLayout{57}{Alarmzeichen {\color{ubuntured}\bfseries{\VarFlag ~}} --
Übersicht}{}{} {
\begin{tabularx}{\linewidth}{XXX}
\gTabHeading{Massive Störung}
&
\gTabHeading{Typische Symptome}
&
\gTabHeading{Gefährliche Diagnosen}
\\\midrule
Bewusstsein

Atemweg / Atmung 

Kreislauf 
&
Atemnot, die sich nicht bessert 

Brodelndes Atemger\"ausch 

Thoraxschmerz  

Schocksymptome 

entgleiste Vitalwerte

schwere Verletzungen

{\dots} 
&
Herzinfarkt

Kardiales Lungenödem

Status Epilepticus

Beckenfraktur

\ldots
\\\bottomrule
\end{tabularx}
}{{\color{ubuntured}\bfseries{\VarFlag}} \gCa{Alarmzeichen:} Wann ist eine vitale Bedrohung
wahrscheinlich? Eine Übersicht.\label{tab:alarmzeichen}}{}




W\"ahrend des weiteren Verlaufs (Anamnesegespr\"ach, etc.) kann man auf
weitere Befunde stoßen welche auf eine vitale Be\-drohung hinweisen
(zB. massiv erhöhter Blutdruck bei vor\-ge\-sch\"adigtem Herzen,
plötzliche oder schleichende Zustands\-verschlechterung, etc.). 

Die Ersteinsch\"atzung muss st\"andig überprüft werden. Neue Befunde
und der Zustand des Patienten können die Situation \"andern.

\subsubsection{Zweiter Schritt: Erstmaßnahmen}
\label{topic:patientenmanagement-erstmassnahmen}

Welche Maßnahmen sind für den Patienten am wichtigsten und
müssen zuerst erfolgen? 

\marginpar{\gH{Bei der Prüfung "`blind"' auf
Nummer sicher gehen und alle gleich für "`vital bedroht"'
erkl\"aren gilt nicht!}

Viele Patienten bei der Prüfung sind nicht vital bedroht und brauchen
keinen Notarzt. 

Erkl\"art man einen Patienten für vital bedroht muss man das
begründen können!
}


Bei vital bedrohten Patienten gibt es 
\begin{center}
\gEs{5 \gT{\gTxMassVit}},
\end{center}
die sobald als möglich (evt. sogar
noch vor Abschluß des Ersteinschätzungsblockes) zu erfolgen haben. 


\gMassnahmenNG{Standardmaßnahmen}{\gTxMassVit}{}{
\label{tab:standardmassnahmen-vital}\label{topic:standardmassnahmen-vital}
\begin{center}

\begin{tabulary}{\linewidth}{LL}
\toprule
\multicolumn{2}{c}{\gTabHeading{\gTxMassVit}}
\\\midrule
Situationsgerechte Lagerung
 &
\\\addlinespace
Sauerstoffgabe
 &
je nach Indikation, Vorsicht bei COPD und Hyper\-ventilation
\\\addlinespace
Notarzt{}-Nachforderung
 &
mit kurzer Begründung
\\\addlinespace
Reanimationsbereitschaft
 &
Platz schaffen, Beatmungsbeutel zu\-sammen\-bauen, ggfs. Defi und
Absaugung holen und in Griffweite stellen
\\\addlinespace
Vitalwerte erheben und Monitoring
 &
RR, HF, Pulsoyxymetrie wenn vor\-handen
\\\bottomrule
\end{tabulary}
\end{center}
}

Bei
nicht vital bedrohten Patienten muss individuell überlegt werden welche
Maßnahmen zuerst zu treffen sind. 

\subsubsection{Dritter Schritt: Durchatmen, nachdenken,
planen}

\emph{Was ist wichtig, was kann warten?} -- Weitere Untersuchungen, S-KAMEL und
Maßnahmen.

Nach den Erstmaßnahmen muss man nun planen, welche Maßnahmen,
Untersuchungen und Fragen am dringlichsten sind und zuerst durchgeführt
werden müssen, und was eher in der Abfolge nach hinten ver\-schoben
wird. 

Die folgende Checkliste in Tabelle
\ref{tab:checkliste-patientenmanagement} gibt einen
Überblick über die Maßnahmen und das Anamnese\-ge\-spr\"ach, die Reihenfolge ist bei jedem Patienten anders!

%%%%%%%%%%%%%%%%%%%%%%%%%%%%%%%%%%%%%%%%%%%%%%%%%%%%%%%%%%%

\begin{gWideParEnv}
\begin{minipage}{\linewidth}
\gCaptionof{table}{Checkliste Patientenmanagement --
Mögliche Maßnahmen, die Reihenfolge ist je nach Patient unterschiedlich\label{tab:checkliste-patientenmanagement}}

{\sffamily
\begin{tabulary}{\linewidth}[H]{CCL}
\toprule\addlinespace
\gEs{Selbstschutz}
%\includegraphics[width=1.005cm,height=0.9cm]{AASdev20090228-img14.gif}
 &
 &
Schutzausrüstung (Handschuhe, Uniform, Helm)
Gefahrenbereich, Unfallstelle absichern, Hunde wegsperren lassen
ggfs. Rückzug, Spezialkr\"afte anfordern, {\dots}
{\dots}
\\\addlinespace\midrule\addlinespace
\gEs{Erkennen}

\includegraphics[width=1cm]{icons/eye.png}
 &
 &
Situation/Milieu {\tiny(zB Alter, soziales Umfeld, SMV,
Alkohol, Pensionistenheim, ...)}

offensichtliche Symptome {\tiny(zB Zyanose,
Bewusstlosigkeit, starke Blutung, ...)}

Schockgefahr {\tiny(Zyanose, kalter Schweiß,
Radialispuls tastbar/nicht tastbar, Puls beschleunigt)}
\\\addlinespace\midrule\addlinespace
\gEs{Erste Hilfe}

\includegraphics[width=1cm]{icons/firstAid.png}
 &
 &
Lagerung

Helmabnahme

Blutstillung

psychischer Beistand

Wärmeerhaltung / Kühlung
\\\addlinespace\midrule\addlinespace
{\multirow{6}{2cm}{\centering\includegraphics[width=1cm]{icons/exclamationMark.png}
\gEs{Anamnese}}}
 &
{\sffamily\textbf S}
 &
Symptome \& Schmerzen 

{\tiny(Zeit/Verlauf, St\"arke, Art/Qualit\"at,
Lokalisation/Ausstrahlung, Provokation)}
\\\addlinespace
 &
{\sffamily\textbf K}
 &
Krankheiten {\tiny(Bestehen Grunderkrankungen, zB
Diabetes?)}
\\\addlinespace
 &
{\sffamily\textbf A}
 &
Allergien/Allergische Reaktion (auch
Allergien auf Medikamente)
\\\addlinespace
 &
 {\sffamily\textbf M}
 &
Medikamente {\tiny(Nimmt Patient Medikamente?
Wogegen sind diese?)}
\\\addlinespace
 &
 {\sffamily\textbf E}
 &
Ereignisse vor Notfalleintritt {\tiny(zB
Anstrengung vor Brustschmerz, Sturz, ...)}
\\\addlinespace
 &
 {\sffamily\textbf L}
 &
Letzte ... (zB letzte Mahlzeit, letzte
Regelblutung, letzter Spitalsaufenthalt)
\\\addlinespace\midrule\addlinespace

{\multirow{8}{2cm}{\centering\includegraphics[width=1cm]{icons/stethoskop.jpg}
\gEs{Befunde}}}
 &
{\sffamily\textbf RR}
 &
auskultatorische und palpatorische Messung
\\\addlinespace
 &
{\sffamily\textbf HF}
 &
Herzfrequenz (Puls), Puls rhythmisch oder arrhytmisch
\\\addlinespace
 &
{\sffamily AF}
 &
Atemfrequenz ausz\"ahlen
\\\addlinespace
 &
{\sffamily Sp\ce{O2}}
 &
Sauerstoffs\"attigung im Blut, Monitoring
\\\addlinespace
 &
{\sffamily  BZ}
 &
Blutzucker {\tiny(Wann war die letzte
Mahlzeit/Insulinspritze? Diabetes bekannt?)}
\\\addlinespace
 &
{\sffamily Temp.}
 &
Körpertemperatur/Fieber
\\\addlinespace
 &
 {\sffamily\small Neuro\-status}
 &
Bewusstsein (ev. GCS), Orientierung,
Pupillen, Halbseitenzeichen, Herdblick, Meningismus, retrograde Amnesie
\\\addlinespace
 &
{\sffamily\small Trauma\-check}
 &
abnorme Beweglichkeit, DMS an allen
Extremit\"aten (Nagelbettprobe), Wunden, paradoxe Atmung,
Abwehrspannung des Bauches, ...
\\\addlinespace\midrule\addlinespace
 \gEs{Sanitäts\-hilfe}

\includegraphics[width=1cm]{icons/san.png}
&
&
Entscheidung Notarzt-Nachforderung

Fortführung der Ersten Hilfe

Schienung

Sauerstoffgabe

Reanimationsbereitschaft

Verlaufskontrolle

\ldots
\\\addlinespace\midrule\addlinespace
\gEs{Notarzt\-assistenz}
\includegraphics[width=1cm]{icons/asb.png}
&
 &
Infusion/periphervenösen Zugang vorbereiten

Intubation vorbereiten

Medikamente vorbereiten

EKG anlegen

\ldots
\\\addlinespace\midrule\addlinespace
\gEs{Krankenhaus}

\includegraphics[width=1cm]{icons/hospitalStaff.png}
 &
 &
Welches Spital/Bett buche ich ab? (Arbeitsunfall? Welche Abteilung
fahre ich an? Ist der Patient bereits in einem Spital in Behandlung?)

Brauche ich eine Spezialabteilung? (Schockraum, Stroke Unit, PTCA,
Verbrennungsstation)

Arztübergabe (Anamnese, Befunde, Maßnahmen, Hinweise)
\\\addlinespace\bottomrule
\end{tabulary}}
\end{minipage}
\end{gWideParEnv}

\cite{AMLS3}, \cite{PHTLS6}, \cite{emergency9}, \cite{SemmelABC1}


% \subsection{Der Patient endet nicht mit der Übergabe}
% 
% \begin{quote}
% Aortenaneurysma
% \end{quote}
% 
% \begin{quote}
% Aortenaneurysma 2
% \end{quote}
% 
% \subsection{Fehler und Fallen}
% vdfvrwbwrtbwrtbrwtb
% 
% \begin{quote}
% Ketanest --- {\itshape Eines Tages, es war Winter, die Straßen verschneit und
% glatt, wurden wir zu einer Patientin mit einer offenen Unterschenkelfraktur
% berufen. Die Dame war auf einer nassen Stufe im Hauseingang augerutscht. Sie
% wirkte soweit gefasst, aufgrund des Schocks\footnote{Schock: hier
% psychologischer Schock.} schien sie noch keine großen Schmerzen wahrzunehemn.
% Trotzdem wurde zur prophylaktischen Schmerztherapie ein Notarzt nachbestellt.
% Dieser entschloß sich zu einer Sedierung, u.a. mit
% Dormicum{\texttrademark}\footnote{Dormicum{\texttrademark}, Wikstoff Midazolam,
% ist ein Schlafmittel, welches zur Sedierung und Narkose verwendet wird.}. (An
% die anderen Medikemnte kann ich mich nicht mehr genau erinnern, ich denke
% Ketanest{\texttrademark}\footnote{Ketanest{\texttrademark}, Wikstoff S-Ketamin,
% ist ein Narkosemittel.} war auch mit ihm Spiel.) 
% 
% Ich erninnere mich noch dass der Notarzt etwas vom Dormicum{\texttrademark}
% spritzte und sich nach einiger Zeit wunderte weildie erwünschte Wirkung
% ausblieb. Er spritzte noch etwas nach und langsam begann es zu wirken \ldots 
% \emph{"`Na bumm, die hat aber viel gebraucht"'}, merkte er nach der Übergabe
% dann  noch an. Des Rätsels Lösung war aber nicht bei der Patientin zu suchen,
% sondern bei den verwendeten Ampullen: 
% 
% Wir vewendeten Ampullen mit 5\,mg Wirkstoff, die Organisation, der der Notarzt
% angehörte, Ampullen mit 15\,mg (!) Wirkstoff. Daher war es wenig verwunderlich,
% dass "`unser"' Dormicum{\texttrademark} viel weniger Wirkung hatte. Nur gut
% dass der Notarzt "`sein"' Ketanest{\texttrademark} aus der Notarzt-Tasche
% verwendet hatte: Unsere Ketanest{\texttrademark}-Ampullen hatten einen Inhalt von 50\,mg, die des
% Notarztes 25\,mg {\ldots} }
% \end{quote}


\clearpage


\subsubsection{Ein Beispiel \ldots}
\paragraph{Ein Beispiel: Herr Sepp hat Atemnot}
\gMargin{\includegraphics[width=\linewidth]{./images/teaser/rtw-asb-921-alt.jpg}}
\begin{verbatim}
Zeit: dd.04.2003 14:15
Code: 26-B-1 
   Kranke Person (Differentialdiagnose): 
   Unbekannter Zustand (Anrufer 3. Hand) 
BO  : 1210, xxxstraße 13/8/8/45, WHG

\end{verbatim}

\begin{quote}
\marginlabel{Szeneüberblick}Sie kommen in eine Wohnung. Die Umgebung erscheint
sicher. Die Wohnung sieht sauber und ordentlich aus. Auf der Couch liegt flach ein
\marginlabel{Ersteindruck}50-jähriger Mann mit weit aufgerissenen Augen und
sichtlicher Atemnot. Er wirkt blass und schweißig und fasst sich mit der
rechten Hand an den Brustbereich.

Sie begrüßen den Patienten, stellen sich vor und
\marginlabel{Hauptbeschwerde}fragen welche Beschwerden er hat. Der Patient
\marginlabel{Atemweg, Atmung}atmet schnell und ringt nach Luft, besondere
Atemgeräusche sind nicht hörbar. Er klagt über Schmerzen in der Brust. Diese
haben vor einer Viertelstunde angefangen. \marginlabel{Kreislauf}Währenddessen
haben Sie an der Radialarterie des Patienten den Puls getastet und eine Rekap-Probe gemacht. Der Puls ist sehr schwach, schnell und unregelmäßig. Die Rekap-Zeit war etwas verlangsamt.
\end{quote}

Welche Informationen haben Sie bis jetzt gewonnen?
\begin{quote}
Der Patient ist bei Bewusstsein. Die Atemwege scheinen frei, aber er hat
Atemnot und einen plötzlich aufgetretenen Thoraxschmerz. Er ist blass und
schweißig, vielleicht steckt ein Schockgeschehen dahinter. Der schwach fühlbare
Puls und die verlängerte Rekap-Zeit unterstützen diese Vermutung.
\end{quote}

Ist der Patient vital bedroht?
\begin{quote}
\marginlabel{Vitale Bedrohung: frühzeitig erkannt (noch vor der "`schnellen
Untersuchung"')}Ja. Zwar wurden noch keine Vitalwerte erhoben, aber bereits
jetzt gilt der Patient als vital bedroht aufgrund der plötzlich einsetzenden Thoraxschmerzen und der Atemnot. Außerdem lässt die Kaltschweißigkeit, der
schwache Puls und die verlängerte Rekap-Zeit auch nichts gutes vermuten.
\end{quote}

Können Sie schon etwas tun?
\begin{quote}
\marginlabel{\gTxMassVit}Ja.  Bei erhaltenem Bewusstsein, Atemnot und
Thoraxschmerz wird der Patient mit erhöhtem Oberkörper gelagert. Weiters sprechen die
Symptome für den Einsatz von Sauerstoff. Es muss noch schnell
geklärt werden ob eine Hyperventilationstetanie (nicht
wahrscheinlich) oder COPD-Erkrankung (erfragen) vorliegt.

\gTxMassVit:

\begin{itemize}
    \item Der Oberkörper des Patienten wird hochgelagert (Atemnot,
    Thoraxschmerz)
    \item Der Notarzt wird nachgefordert (Berufungsgrund: akuter
    Thoraxschmerz mit Atemnot und schlechtem AZ).
    \item Der Patient wird gefragt ob er an COPD leidet, dies wird verneint.
    Eine psychisch bedingte Hyperventilation erscheint unwahrscheinlich.
    
    Daher wird eine Sauerstoffmaske mit einem Flow von 10l/min angelegt.
    \item Reanimationsbereitschaft wird hergestellt
    \item Vitalwerte: siehe nächster Schritt.
\end{itemize} 

\end{quote}


Was sind die n\"achsten Schritte?
\begin{quote}
\marginlabel{Schnelle Untersuchung, Nicht-Traumapatient: Vitalwerte}Wichtig sind
nun die Vitalwerte und das frühe Beginnen des Monitorings.Dann muss man genaueres über das Beschwerdebild in
Erfahrung bringen. Der weitere Plan h\"angt davon ab.

Die Vitalwerte werden gemessen und ein Pulsoxy angeh\"angt (RR 100/60,
HF 110, SpO2 94\%). Der Patient gibt an, dass er st\"arkste
stechende/brennende Schmerzen hinter dem Brustbein hat die ins Kiefer
ausstrahlen. Außerdem fühle es sich so an als würde jemand auf dem
Patienten drauf stehen. Die Schmerzen sind nicht atem- oder
bewegungsabh\"angig  und haben vor einer Viertelstunde
{\quotedblbase}wie aus dem Nichts heraus{\textquotedblleft} begonnen.
Außerdem bekomme er schlecht Luft, sonst h\"atte er keine Beschwerden.
\end{quote}

\marginlabel{Problem des Patienten}Als Leitsymptom erkennen Sie den
Thoraxschmerz und die Atemnot.

An welche Differentialdiagnosen müssen Sie nun denken?


% \clearpage
% CHECK gWideParEnv Thoraxschmerz: Differentialdiagnosen unseres Patienten
\begin{gWideParEnv}
\begin{tabularx}{\linewidth}[H]{XXXc}
\toprule
\multicolumn{4}{c}{\centering\gTabHeading{Thoraxschmerz: Differentialdiagnosen unseres
Patienten} }
\\\midrule
 &
\centering \gTabSub{dafür spricht}
 &
\centering \gTabSub{dagegen spricht}
 &
\gTabSub{Bewertung}
\\\midrule
\emph{Herzinfarkt, Angina pectoris}
 &
Schmerzen hinter dem Brustbein, Ausstrahlung, nicht atemabh\"angig,
Druck-/Engegefühl, Atemnot
 &
 &
sehr wahrscheinlich
\\\midrule
\emph{Lungenembolie}
 &
Schmerzen hinter dem Brustbein, Atemnot
 &
untypische Ausstrahlung, 
 &
gut möglich
\\\midrule
\emph{Pneumothorax}
 &
Schmerzen im Brustbereich, Atemnot
 &
plötzliche stärkste Schmerzen
 &
eher unwahrscheinlich
\\\midrule
\emph{Trauma}
 &
Schmerzen
 &
kein Hinweis
 &
unwahrscheinlich
\\\midrule
\emph{Erkrankung des Stütz- und Bewegungsapparats}
 &
Schmerzen 
 &
untypische Ausstrahlung, Atemnot, nicht bewegungsabhängig
 &
sehr unwahrscheinlich
\\\bottomrule
\end{tabularx}
\end{gWideParEnv}


\gFazit{Die wahrscheinlichsten Diagnosen sind der
Herzinfarkt und die Lungenembolie.}

Fallen Ihnen bei diesen Erkrankungen wichtige weitere Untersuchungen
ein?
\begin{quote}
\marginlabel{Durchatmen, nachdenken, planen}Bei beiden Erkrankungen ist es
wichtig beide Beine zu begutachten. Im Falle einer Lungenembolie sucht man nach einer Beinvenenthrombose
(einseitig überwärmtes, bläulich-rötliches geschwollenes Bein).
Bei einer Herzinsuffizienz bzw. Herzinfarkt beurteilt man den
Rückstau des Blutes, dh. ob die Beine geschwollen sind.
\end{quote}


Außerdem, hat der Patient so etwas schon einmal gehabt? Oder sind einschlägige
Erkrankungen bekannt?  

\begin{quote}
Sie begutachten beide Beine des Patienten und stellen mittels
Daumendruckprobe fest, dass beide Beine leicht geschwollen sind. Auf
Nachfrage gibt der Patient an, dass er sonst keine geschwollenen Beine
hat. 
\end{quote}

Was sind die nächsten Schritte?

\begin{quote}
Die wichtigsten Untersuchungen wurden gemacht, nun muss man sich um die
Anamnese kümmern. 

Symptome \& Schmerzen: siehe oben.

Krankheiten: Der Patient leidet an nicht-Insulinpflichtigem Diabetes
mellitus und Bluthochdruck. Sonst sind keine Krankheiten bekannt.
Spitalsbriefe gibt es nicht. 

Sie merken sich, dass Sie nachher den Blutzucker messen wollen!

Allergien: keine bekannt.

Medikamente: Glucophage 850\,mg für den Zucker, Tritace für den
Blutdruck. Medikamente wurden heute wie verschrieben eingenommen. 

Ereignisse vor Notfalleintritt: Schmerzen begannen plötzlich als der
Patient Wäsche waschen wollte. In den letzten Wochen immer wieder
Stiche in der Herzgegend, hat er aber nicht ernst genommen. 

Letzte(r):

Spitalsaufenthalt: Vor 10 Jahren in der Rudolfstiftung wegen Rotlauf.

Arztbesuch: Vor einer Woche wegen Neuausstellung eines Rezeptes.

Mahlzeit: Mittagessen, Gulasch

Regel: Der Patient ist männlich!
\end{quote}
% pro:
% 
% Lungenembolie 

Was glauben Sie wie es weiter geht?

\gTakeNotesLarge

% ========Subsection eingefügt vom Emhofer 1.5.2009===============================

\section{Übergabegespräch an weiterbehandelndes Fachpersonal}  

\begin{wrapfigure}{r}{5cm}
\includegraphics[width=\linewidth]{./images/shared/aufnahmearzt-ihres-vertrauens-edited2-0500.jpg}
\end{wrapfigure}Die Übergabe des Patienten in den Verantwortungsbereich eines Notarztes oder des
Spitalspersonals hat geordnet und gewissenhaft zu erfolgen, damit keine wichtige
Information verloren geht. Diese Schnittstelle zwischen Sanitäter und Arzt bzw.
Krankenschwester kann auf den weiteren Behandlungsverlauf einen großen Einfluss
nehmen und darf daher in seiner Bedeutung nicht unterschätzt werden! Es empfiehlt
sich -- so Zeit dafür ist -- sich bereits vor Eintreffen des Notarztes oder vor
Ankunft im Krankenhaus zu überlegen, welche Informationen wesentlich sind und
daher auf keinen Fall vergessen werden dürfen.

Das Übergabegespräch muss mindestens folgende Informationen beinhalten:
\begin{itemize}
    \item Name des Patienten
    \item Alter
    \item Zusammenfassung der wesentlichen Inhalte der Anamnese
    \item Beschreibung der Symptomatik und deren Verlauf
    \item Bisher erhobene Befunde (Blutdruck, Puls, Sauerstoffsättigung usw.)
    \item Bisher durchgeführte Maßnahmen
    \item Sonstige, wichtige Hinweise (z.B. Infektionsgefahren,
    Bluterkrankheit, Allergien etc.)
\end{itemize}

%=========Ende Abschnitt Emhofer==========================================================

\section*{Weiterführende Literatur}
\cite{SudoweProfHandelnRD1,
AMLS3,
PHTLS6,
emergency9,
SemmelABC1,
KirschnickPflege2,
RueckerBildatlas1}